\documentclass[a4paper,12pt]{article}
\usepackage[T2A]{fontenc}
\usepackage[utf8]{inputenc}
\usepackage[russian]{babel}
\usepackage{amsmath,amssymb,mathtools}
\usepackage{geometry}
\geometry{left=24mm,right=24mm,top=23mm,bottom=23mm}
\usepackage{microtype}
\usepackage{booktabs,tabularx,array}
\usepackage{enumitem}
\usepackage{fancyhdr}
\usepackage{xcolor}
\usepackage{tikz}
\usetikzlibrary{arrows.meta,positioning,shapes.geometric,calc}
\usepackage[colorlinks=true,linkcolor=tealblue,urlcolor=tealblue]{hyperref}
\definecolor{tealblue}{HTML}{007E82}
\definecolor{tealdark}{HTML}{00565A}
\definecolor{ink}{HTML}{303437}
\definecolor{muted}{HTML}{5E6467}
\color{ink}
\setlength{\parindent}{0pt}
\setlength{\parskip}{7pt plus 2pt minus 1pt}
\setlist[enumerate]{leftmargin=*,itemsep=9pt,topsep=5pt}
\pagestyle{fancy}
\fancyhf{}
\fancyhead[L]{\small Отношения, доли и проценты}
\fancyhead[R]{\small Надежда Клименко}
\fancyfoot[C]{\small\thepage}
\renewcommand{\headrulewidth}{0.45pt}
\renewcommand{\headrule}{\hbox to\headwidth{\color{tealblue}\leaders\hrule height \headrulewidth\hfill}}
\setlength{\headheight}{25pt}
\newcommand{\partsection}[1]{\vspace{10pt}\section*{\textcolor{tealdark}{#1}}\vspace{-2pt}\noindent\textcolor{tealblue}{\rule{\linewidth}{0.5pt}}\vspace{3pt}}
\newcommand{\focus}[1]{\textcolor{tealdark}{#1}}
\newcommand{\ex}[1]{\par\vspace{4pt}{\focus{\textbf{Пример.} #1}}\par}
\newcommand{\ruletext}[1]{\textbf{Главное правило.} #1}
\newcommand{\step}[1]{\textbf{#1}}
\begin{document}
\begin{titlepage}
\thispagestyle{empty}
\begin{center}
\vspace*{23mm}
{\Large\color{tealdark} МАТЕМАТИКА\par}
\vspace{24mm}
{\LARGE Чек-лист\par}
\vspace{19mm}
{\LARGE\color{tealblue} Отношения, доли\par}
\vspace{3mm}
{\LARGE\color{tealblue} и проценты\par}
\vspace{18mm}
{\large Текстовые задачи и сравнение величин\par}
\vfill
{\large Надежда Клименко\par}
\vspace{12mm}
{\large 17 ноября 2026 г.\par}
\vfill
{\normalsize Лёвин Артём Александрович\par}
{\normalsize эксклюзивно для Надежды Клименко\par}
\vspace{18mm}
\begin{tikzpicture}[x=1cm,y=1cm]
 \draw[tealblue,line width=0.75pt] (-6,0) .. controls (-4,0.95) and (-2,-0.75) .. (0,0)
 .. controls (2,0.8) and (4,-0.9) .. (6,0);
 \fill[tealdark] (0,0) circle (1.4pt);
\end{tikzpicture}
\vspace{8mm}
\end{center}
\end{titlepage}
\setcounter{page}{1}

\partsection{1. Отношение: как сравнить две величины}
\ruletext{Отношение величины $A$ к величине $B$ --- это результат деления $A$ на $B$:}
\[
A:B=\frac{A}{B},\qquad B\ne0.
\]
Читаем условие внимательно: слова \emph{«отношение моркови к капусте»} означают, что в числителе стоит масса моркови, а в знаменателе --- масса капусты.

\ex{Для приготовления взяли $0{,}7$ кг моркови и $4{,}9$ кг капусты.}
\[
0{,}7:4{,}9=7:49=1:7.
\]
На одну массовую часть моркови приходится семь таких же частей капусты. При сохранении этого отношения можно взять, например, $1$ и $7$ кг или $2$ и $14$ кг: масштаб меняется, а соотношение остаётся прежним.

\partsection{2. Часть от целого}
\ruletext{Чтобы найти, какую часть целого составляет величина, делим величину части на величину целого:}
\[
\boxed{\text{Доля}=\frac{\text{часть}}{\text{целое}}}.
\]
Часть и целое должны измеряться в одинаковых единицах. Если часть содержится в целом, её доля находится между нулём и единицей включительно.

\ex{Катя выполнила $18$ прыжков, из них $3$ неудачно. Даша выполнила $21$ прыжок, из них $2$ неудачно.}
Удачных прыжков: у Кати $18-3=15$, у Даши $21-2=19$. Их доли успеха:
\[
\text{Катя: }\frac{15}{18}=\frac56,\qquad
\text{Даша: }\frac{19}{21}.
\]
Для сравнения приведём дроби к знаменателю $42$:
\[
\frac56=\frac{35}{42}<\frac{38}{42}=\frac{19}{21}.
\]
\focus{Вывод:} у Даши доля успешных прыжков выше. Сравнение только количества удачных прыжков может вводить в заблуждение, если число попыток различается.

\partsection{3. Отношение частей и масштаб}
Если стороны прямоугольника относятся как $2:1$, их можно обозначить $2x$ и $x$, где $x>0$. При $x=1$ получаем $2$ и $1$ см; при $x=4$ --- $8$ и $4$ см.

\ex{Первая бригада построила $\tfrac59$ тоннеля, вторая --- оставшуюся часть. Во сколько раз первая часть больше второй?}
Вторая бригада построила $1-\tfrac59=\tfrac49$ тоннеля. Формулировка \emph{«во сколько раз»} указывает на деление:
\[
\frac{\tfrac59}{\tfrac49}=\frac59\cdot\frac94=\frac54=1{,}25.
\]
\focus{Итог:} первая бригада построила в $1{,}25$ раза больше второй.

\partsection{4. Процент как способ записи доли}
Один процент --- одна сотая часть: $1\%=\tfrac1{100}$.
\[
\boxed{\text{Процентное содержание}=\frac{\text{часть}}{\text{целое}}\cdot100\%}.
\]
Если известна доля $\tfrac14$, то ей соответствует $25\%$; доле $\tfrac34$ --- $75\%$. Обратный переход: $35\%=\tfrac{35}{100}=\tfrac7{20}$.

\ex{Ученик прочитал в три раза меньше страниц, чем осталось прочитать.}
Удобно обозначить число прочитанных страниц через $x$, оставшихся --- через $3x$. Тогда объём книги равен $4x$, причём $x>0$.
\[
\begin{aligned}
\text{Прочитано:}\quad&\frac{x}{4x}=\frac14=25\%,\\[3pt]
\text{Осталось:}\quad&\frac{3x}{4x}=\frac34=75\%.
\end{aligned}
\]
А какую часть \emph{от прочитанных} страниц составляет оставшееся?
\[
\frac{3x}{x}=3=300\%.
\]
\focus{Важно:} процент может превышать $100\%$, если сравниваемая величина больше выбранной базы. Здесь базой выступают прочитанные страницы, поэтому результат $300\%$ корректен.

\partsection{5. Проценты одного количества относительно другого}
Пусть первый отряд выполнил $45\%$ общего объёма работ, а два других вместе --- $55\%$. Сколько процентов составляет работа первого отряда \emph{от работы двух других}?
\[
\frac{45\%}{55\%}\cdot100\%=\frac9{11}\cdot100\%\approx81{,}8\%.
\]
Сначала делим сравниваемые части (одинаковые обозначения процентов сокращаются), затем выражаем найденное отношение в процентах.

\textbf{Как распознать знаменатель:} слова «от площади двух других», «от числа попыток», «от прочитанных страниц» называют \emph{базу сравнения}. Именно она стоит в знаменателе.

\textbf{Не путайте формулировки:}
\begin{itemize}[leftmargin=1.5em,itemsep=2pt]
\item «Какую часть $A$ составляет от $B$?» --- $\tfrac AB$.
\item «Во сколько раз $A$ больше $B$?» --- $\tfrac AB$ (ответ в разах).
\item «Сколько процентов $A$ составляет от $B$?» --- $\tfrac AB\cdot100\%$.
\end{itemize}

\partsection{6. Массовая доля вещества в растворе}
\ruletext{Масса раствора складывается из масс растворённого вещества и растворителя:}
\[
 m_{\text{раствора}}=m_{\text{вещества}}+m_{\text{воды}}.
\]
\[
\boxed{\omega=\frac{m_{\text{вещества}}}{m_{\text{раствора}}}\cdot100\%}.
\]
\ex{В $1200$ г воды растворили $40$ г соли. Найдите массовую долю соли.}
Масса рассола $1200+40=1240$ г. Следовательно,
\[
\omega=\frac{40}{1240}\cdot100\%=\frac{100}{31}\%\approx3{,}23\%.
\]
\textbf{Проверка здравого смысла:} $40$ г соли составляют малую часть от $1240$ г раствора, поэтому процент должен быть небольшим.

\ex{К тому же рассолу добавили ещё $800$ г воды.}
Масса соли остаётся $40$ г; общая масса становится $2040$ г. Новая массовая доля:
\[
\omega_2=\frac{40}{2040}\cdot100\%=\frac{100}{51}\%\approx1{,}96\%.
\]
При добавлении чистой воды массовая доля соли уменьшается, поскольку масса соли постоянна, а масса раствора увеличивается.

\partsection{7. На сколько процентов величина увеличилась?}
\ruletext{Изменение в процентах вычисляется относительно исходного значения $B$:}
\[
\boxed{\text{Прирост}=\frac{A-B}{B}\cdot100\%,\qquad B>0.}
\]
\ex{В этом году школьники собрали $720$ кг макулатуры --- на $100$ кг больше, чем в прошлом.}
В прошлом году было $720-100=620$ кг. Процент прироста:
\[
\frac{720-620}{620}\cdot100\%=\frac{100}{620}\cdot100\%\approx16{,}1\%.
\]
\focus{Итог:} сбор вырос примерно на $16{,}1\%$ по отношению к прошлому году.

\textbf{Самопроверка:} если требуется «на сколько процентов больше», сначала найдите разность, затем разделите её на \emph{первоначальную} величину и умножьте на $100\%$.

\partsection{8. Частые ошибки}
\begin{enumerate}[label=\textcolor{tealdark}{\arabic*.},itemsep=4pt]
\item Меняют местами числитель и знаменатель. Проверяйте, \emph{что} сравниваем и \emph{с чем}.
\item Сравнивают абсолютные количества при разных общих объёмах. Для честного сравнения используют доли.
\item Приписывают знак процента к результату деления без умножения на $100$. Например, $\tfrac14=0{,}25=25\%$.
\item В задачах о растворах делят массу соли на массу воды. Делить нужно на массу всего раствора.
\item При вычислении прироста делят на новое значение. За базу следует брать исходную величину.
\item Смешивают единицы измерения или неправильно группируют категории, перечисленные в условии.
\end{enumerate}

\clearpage
\begin{center}
{\Large\color{tealdark} Алгоритм: перевод текста в отношение и проценты}\par
\vspace{13mm}
\begin{tikzpicture}[
 node distance=13mm and 17mm,
 every node/.style={align=center,font=\small},
 startstop/.style={ellipse,draw=tealblue,line width=.7pt,minimum width=39mm,minimum height=10mm},
 action/.style={rounded corners=2pt,draw=tealblue,line width=.65pt,text width=95mm,minimum height=14mm,inner sep=6pt},
 decision/.style={diamond,aspect=2.4,draw=tealblue,line width=.7pt,text width=39mm,inner sep=2pt},
 result/.style={rounded corners=2pt,draw=tealblue,line width=.65pt,text width=51mm,minimum height=18mm,inner sep=4pt},
 arr/.style={-{Stealth[length=2.7mm]},line width=.7pt,draw=tealdark}
]
\node[startstop] (begin) {Прочитайте вопрос};
\node[action,below=of begin] (a) {Найдите сравниваемую величину $A$ и базу $B$: \emph{от чего} или \emph{по отношению к чему} считают};
\node[action,below=of a] (b) {Запишите отношение $\displaystyle\frac AB$; убедитесь, что $B\ne0$ и единицы согласованы};
\node[decision,below=18mm of b] (c) {Ответ нужен\\в процентах?};
\node[result,below left=20mm and 6mm of c] (no) {Нет: сократите дробь или запишите, во сколько раз};
\node[result,below right=20mm and 6mm of c] (yes) {Да: умножьте отношение на $100\%$};
\node[startstop,below=33mm of c] (end) {Проверьте смысл ответа};
\draw[arr] (begin)--(a);
\draw[arr] (a)--(b);
\draw[arr] (b)--(c);
\draw[arr] (c.south west)--node[above left,font=\footnotesize]{нет}(no.north east);
\draw[arr] (c.south east)--node[above right,font=\footnotesize]{да}(yes.north west);
\draw[arr] (no.south)|-(end.west);
\draw[arr] (yes.south)|-(end.east);
\end{tikzpicture}
\end{center}
\clearpage

\partsection{Тренировочный блок}
\textit{Путь А. 10 заданий: от базовых отношений к многошаговым задачам.}\par
\textbf{Базовый уровень}
\begin{enumerate}[label=\textbf{\arabic*.},itemsep=10pt]
\item Из $24$ бросков в кольцо успешными оказались $18$. Какую часть составляют успешные броски? Сколько это процентов?
\item Для смеси использовали $0{,}9$ кг одного ингредиента и $3{,}6$ кг другого. Найдите отношение массы первого ингредиента к массе второго в простейшем виде.
\item Стороны прямоугольника относятся как $3:2$. Меньшая сторона равна $8$ см. Найдите большую сторону.
\item У первой команды успешными были $14$ попыток из $18$, у второй --- $21$ из $30$. У какой команды доля успеха выше?
\end{enumerate}
\textbf{Средний уровень}
\begin{enumerate}[label=\textbf{\arabic*.},start=5,itemsep=10pt]
\item Одна бригада построила $\tfrac7{12}$ дороги, другая --- остальную часть. Во сколько раз первая бригада построила больше второй?
\item Первый отряд выполнил $36\%$ всей работы, а второй --- $24\%$. Сколько процентов составляет работа первого отряда от работы второго?
\item Прочитанные и непрочитанные страницы книги относятся как $2:3$. Найдите доли и проценты прочитанных и оставшихся страниц. Сколько процентов составляют оставшиеся страницы от прочитанных?
\end{enumerate}
\textbf{Повышенный уровень}
\begin{enumerate}[label=\textbf{\arabic*.},start=8,itemsep=10pt]
\item В $900$ г воды растворили $60$ г соли. Найдите массовую долю соли. Затем добавили $540$ г воды. Какой стала массовая доля соли?
\item В прошлом году собрали $480$ кг макулатуры, в этом --- $600$ кг. На сколько процентов увеличился сбор?
\item В спортивной школе $27\%$ занимаются футболом, $44\%$ --- баскетболом, остальные --- только неигровыми видами спорта. Сколько процентов составляет количество занимающихся неигровыми видами спорта \emph{от количества занимающихся игровыми}? Ответ округлите до десятых.
\end{enumerate}

\clearpage
\partsection{Ответы к тренировочному блоку}
\renewcommand{\arraystretch}{1.5}
\begin{tabularx}{\linewidth}{@{}>{\bfseries}p{12mm}X@{}}
\toprule
№ & Ответ \\
\midrule
1 & $\tfrac34=75\%$.\\
2 & $1:4$.\\
3 & $12$ см.\\
4 & У первой команды: $\tfrac79>\tfrac7{10}$.\\
5 & $\tfrac75=1{,}4$ раза.\\
6 & $150\%$.\\
7 & Прочитано: $\tfrac25=40\%$; осталось: $\tfrac35=60\%$; оставшиеся страницы составляют $150\%$ от прочитанных.\\
8 & Сначала $6{,}25\%$, после разбавления $4\%$.\\
9 & На $25\%$.\\
10 & $\tfrac{29}{71}\cdot100\%\approx40{,}8\%$.\\
\bottomrule
\end{tabularx}
\vspace{14mm}
\textbf{Краткая памятка}\par
\begin{tabularx}{\linewidth}{@{}>{\raggedright\arraybackslash}p{.40\linewidth}>{\raggedright\arraybackslash}X@{}}
\toprule
Сигнал в задаче & Ваше действие \\
\midrule
«Какую часть составляет?» & Часть разделите на целое.\\
«Во сколько раз больше?» & Большую величину разделите на меньшую.\\
«Сколько процентов составляет?» & Отношение умножьте на $100\%$.\\
«Сколько процентов соли?» & Массу соли разделите на массу раствора и умножьте на $100\%$.\\
«На сколько процентов выросло?» & Прирост разделите на исходное значение и умножьте на $100\%$.\\
\bottomrule
\end{tabularx}
\vfill
\begin{center}\small\color{muted}Лёвин Артём Александрович --- эксклюзивно для Надежды Клименко\end{center}
\end{document}
